% Noether P09 Korean producer draft — UNCHECKED
% T01-U03; ED0002 lines6370--6383; source 1,011 LF B / AE94C006A3F137D23EC76002A6E9A3C41DC9DF9DA47849DDA41FE9BF09B1B6EE
% Pointer v009 / ED0002: B06BE3530D9CF2E82B56FDBA7FE41D5D044DF2425DFA2A059D4939EAA2F7A6C2 / C9A125167ACB33D914EE4374B65AE7CDF0052F568371B8B77B720EA178ABF0E3
% Translation only; algebraische Basis is provisionally rendered 초월기저 in its local Steinitz sense.

이 구성은 정렬정리로 모든 수의 \emph{초월기저}가 존재함을 알 수 있다는 사실에 바탕을 둔다. 다시 말해, 이는 수 $\eta$들의 \emph{계} $H$로서, \emph{그 수들 사이에는 아무런 대수적 관계도 없지만 그 수들을 써서 다른 모든 수를 대수적으로 나타낼 수 있는} 것이다. 이 기저의 수 $\eta$들은 대수적으로 독립이므로 \emph{부정원}으로 볼 수 있고, 따라서 찾는 영역은 계수가 모든 대수적 수의 체 $K$에 속하는, 부정원 $\eta$의 \emph{유리함수와 대수함수로 이루어진 정역}이 된다.\srcfn{**)}{이 때문에 본 논문의 고찰은 저자의 앞선 논문 「Körper und Systeme rationaler Funktionen」, Math. Ann. 76, 161쪽 (1915)의 고찰과 부분적으로 나란히 진행된다.} 그러면 체르멜로가 구성한 특수한 영역 $\mathfrak G_\eta$는 다음 \emph{기저 성질}들로 특징지어진다.
